% !PRINTED 474
\weberpage{474}
%<<CONT
verifiable; but finding the factors is also easy, at least as far as the calculations
have so far been carried, since in later cases one can often use results found earlier
and moreover knows the general form of the factors. Examples will make this clear.

\textbf{§\,128. The first cases of the calculation of $f(\sqrt{-m})$.}

If in the cubic equation:
\weq{1}{u^3 - \gamma_2(\omega)\,u - 16 = 0,}
whose roots according to § 54 are
\[
f(\omega)^8, \quad -f_1(\omega)^8, \quad -f_2(\omega)^8
\]
we put $\omega = i$, hence [§\,125, (16)], $\gamma_2 = 12$, then it follows:
\[
u^3 - 12u - 16 = 0,
\]
an equation whose only positive root is $u = 4$, so that
in agreement with § 126 we obtain:
\weq{2}{f(i) = \sqrt[4]{2};}
if we put $\omega = \sqrt{-2}$, hence $\gamma_2 = 20$ [§\,125, (17)], then from (1) we obtain
the equation:
\[
u^3 - 20u - 16 = 0
\]
with the only rational root $-4$. But since, by the
theorems of the two preceding sections, $f_1(\sqrt{-2})^8$ must be
rational, we have:
\weq{3}{f_1(\sqrt{-2}) = \sqrt[4]{2}.}

We further put $\omega = e^{\frac{2\pi i}{3}}$, hence $\gamma_2(\omega) = 0$ (§ 117), and
obtain from (1):
\[
u^3 = 16.
\]
This equation is satisfied by [§\,34, (19)]:
\[
u = -f_2\left(\frac{-1+\sqrt{-3}}{2}\right)^8 = \frac{16\,e^{\frac{2\pi i}{3}}}{f(\sqrt{-3})^8},
\]
from which the real positive value:
\weq{4}{f(\sqrt{-3}) = \sqrt[3]{2}.}

In order to be able to apply the remaining results of § 125,
we put
\[
m = 7, \quad 11, \quad 19, \quad 43, \quad 67, \quad 163,
\]
%CONT>>
